\section{Transporting excitation tails}
\label{sec:transport}

The next lemma compares two regular trajectories joined at their final
endpoint by a product projection.  Its proof uses the two filters of
Section~\ref{sec:filters} in different directions.  A tail filter $T$
is propagated forward along the projected trajectory, producing an
effective range on which the final projection can be approximated by a
polynomial $D$.  That polynomial is then propagated backwards along the
original trajectory.  The resulting scalar product is small because
$T$ suppresses low counts while the original vector has small high-count
tails.  Positivity of $T$ converts this estimate into a tail bound for
the projected vector.

\begin{lemma}[Tail transport]
\label{lem:tail-transport}
Let $V$ have a bounded number of analytic layers, let $A$ and $B$ be
tests at its initial and final endpoints, respectively, and put
\[
 \chi=V^\dagger P_BV\psi,
 \qquad \norm{\psi}\le1.
\]
Assume that both the trajectory from $\psi$ and the trajectory from
$\chi$ are regular for $V$.  Fix $0\le b<1$.  Suppose that, for every
fixed $b<h<1$ and $0<p<h$, eventually
\begin{equation}
 \label{eq:input-tail}
 \norm{\one_{H_A\ge m^h}\psi}\le\exp(-m^p).
\end{equation}
Then, for every such fixed pair $h,p$, eventually
\[
 \norm{\one_{H_A\ge m^h}\chi}\le\exp(-m^p).
\]
The tests and vectors may vary along the sequence, and $\chi$ is not
normalized.
\end{lemma}

Figure~\ref{fig:transport} records the two uses of regularity.  The
forward propagation concerns $T\chi$, whereas the reverse propagation
concerns $DV\psi$.  Neither filtered vector is assumed to have a regular
trajectory.
In the diagram, $K$ is the degree exponent produced by forward propagation.
The zero filter $D$ has degree at most $m^r$.
Reverse propagation approximates $V^\dagger DV\psi$ by $G\psi$, where $G$
has degree at most $m^{r'}$; the proof below constructs these objects and
proves the displayed bounds.

\begin{figure}[tbp]
\centering
\begin{tikzpicture}[
  >=Stealth,
  every node/.style={font=\small},
  state/.style={draw,rounded corners=2pt,minimum width=3.4cm,
    minimum height=1.15cm,align=center,inner sep=5pt},
  forward/.style={state,draw=blue!65!black,fill=blue!3},
  backward/.style={state,draw=orange!70!black,fill=orange!4}
]
\node[font=\small\scshape] at (0,1.1) {Start};
\node[font=\small\scshape] at (8.5,1.1) {Finish};
\node[forward] (tstart) at (0,0)
  {$T\chi$\\[2pt] $\norm{T\chi}\le 2$};
\node[forward] (tfinish) at (8.5,0)
  {$VT\chi$\\[2pt] effective $H_B$-count range $m^K$};
\draw[->,thick,blue!65!black] (tstart) --
  node[above] {$V$}
  node[below,align=center] {propagate $T$ along $\chi$}
  (tfinish);
\node[backward] (dstart) at (0,-2.2)
  {$G\psi$\\[2pt] $\deg G\le m^{r'}$};
\node[backward] (dfinish) at (8.5,-2.2)
  {$DV\psi$\\[2pt] $D=D(H_B)$: zero filter};
\draw[->,thick,orange!70!black] (dfinish) --
  node[above] {$V^\dagger$}
  node[below,align=center] {vector error $\le\exp(-m^r)$}
  (dstart);
\end{tikzpicture}
\caption{The two transports in Lemma~\ref{lem:tail-transport}.
The forward step establishes a tail bound for $VT\chi$ outside
the count range $m^K$; the projection approximation is estimated
on this vector. The backward step acts on the unfiltered
reference $V\psi$. Its error is paired with $T\chi$, whose norm
is bounded by two.}
\label{fig:transport}
\end{figure}


The proof descends from the automatic zero tail above $m$ to the requested
threshold. At one step, a known tail bound above $m^u$ controls the tail
filter's global growth while the filter detects the new band
$m^h\le H_A\le m^u$. The following lemma supplies compatible exponent gaps
for this step.
Here $u,h$ are the exponents for the old and new count thresholds,
with tail precision exponents $u-\alpha$ and $h-\beta$.
The tail filter has suppression exponent $s$ and degree and logarithmic
norm bounds with exponent $k$, enlarged to $K$ under forward propagation;
the zero filter has precision exponent $q$ and degree and logarithmic
norm bounds with exponent $r$, enlarged to $r'$ under reverse propagation.
Finally, $v_0<v$ are initial count cutoff exponents, and $\theta$ is
the precision exponent used for the assumed input tail above $m^{v_0}$.

\begin{samepage}
\begin{lemma}[Exponent budget]
\label{lem:exponent-budget}
Suppose
\[
 0\le b<h<u\le1,
 \qquad 0<\alpha<\beta<h,
 \qquad u-h<\alpha.
\]
There are fixed exponents $q,k,K,r,r',s,v_0,v,\theta$ satisfying
\begin{gather*}
 0<h-\beta<q<r<r'<s<k<h,
 \qquad k<K<1,
 \qquad k<u-\alpha,\\
 \frac{K+q}{2}<r,
 \qquad s+\frac{u-h}{2}<k,\\
 \max\{b,r'\}<v_0<v<h,
 \qquad r'<\theta<v_0.
\end{gather*}
All these inequalities have strictly positive margins independent of
$m$.
\end{lemma}
\end{samepage}

\begin{proof}
Write $\eta=u-h$ and choose
$\alpha<\lambda<\nu<\beta$.  Set
\[
 k=u-\lambda,
 \qquad q=h-\nu.
\]
Then $q>h-\beta>0$, $k-q=\eta+\nu-\lambda>0$, and
$k<h$ because $\lambda>\alpha>\eta$.  In particular $k<1$.
Choose
\[
 0<e<\min\{\nu-\lambda,1-k\},
 \qquad K=k+e.
\]
The interval available for the remaining filter exponents has width
\[
 w=\left(k-\frac{\eta}{2}\right)-\frac{K+q}{2}
   =\frac{\nu-\lambda-e}{2}>0.
\]
With $a_0=(K+q)/2$, take
\[
 r=a_0+\frac w4,
 \qquad r'=a_0+\frac w2,
 \qquad s=a_0+\frac{3w}{4}.
\]
Thus
\[
 r-a_0=r'-r=s-r'=k-s-\frac\eta2=\frac w4>0.
\]
These choices give $q<r<r'<s<k<h$, as well as
$k<u-\alpha$.  Since $\max\{b,r'\}<h$, choose $v_0,v$
strictly between these endpoints in the stated order and then choose
$r'<\theta<v_0$.  Each choice is made among fixed real numbers.
\end{proof}

\begin{proof}[Proof of Lemma~\ref{lem:tail-transport}]
Fix target exponents $b<h_*<1$ and $0<p_*<h_*$.  Set
\[
 \delta=\frac{h_*-p_*}{4},
 \qquad
 N=\left\lceil\frac{2(1-h_*)}{\delta}\right\rceil,
 \qquad \eta=\frac{1-h_*}{N}.
\]
Here $N\ge1$ is fixed and $0<\eta\le\delta/2$.  Define
\[
 h_i=1-i\eta,
 \qquad \alpha_i=\delta+\frac{i\delta}{N}
 \qquad (0\le i\le N).
\]
We descend from $h_0=1$ to $h_N=h_*$, proving at step $i\ge1$ that
\[
 \norm{\one_{H_A\ge m^{h_i}}\chi}
 \le\exp\bigl(-m^{h_i-\alpha_i}\bigr).
\]
To start, the tail \emph{strictly above} $m^{h_0}=m$ is zero,
since the spectrum of $H_A$ is contained in $\{0,\ldots,m\}$.
No estimate on the eigenspace of count $m$, and no regularity condition
at exponent one, is needed.

Fix a step and abbreviate
\[
 u=h_{i-1},\qquad h=h_i,\qquad
 \alpha=\alpha_{i-1},\qquad \beta=\alpha_i.
\]
Its induction hypothesis is
\[
 \norm{\one_{H_A>m^u}\chi}
 \le\exp\bigl(-m^{u-\alpha}\bigr).
\]
For $i>1$ this follows from the preceding, stronger inclusive bound.
The parameters obey the hypotheses of Lemma~\ref{lem:exponent-budget}:
indeed, $u-h=\eta<\delta\le\alpha$ and
$\beta\le2\delta<h_*\le h$.  Choose its exponents once for this
step.  All subsequent estimates are for sufficiently large $m$.

\emph{Forward propagation on the trajectory of $\chi$.}
Apply Lemma~\ref{lem:tail-filter} with parameters $v,h,u,s,k$, and
write $T=T(H_A)$ for its polynomial evaluated at the initial count.
It is self-adjoint and nonnegative.  Its degree is at most $m^k$,
its global norm is at most $\exp(m^k)$, and
\[
 \norm{\one_{H_A\le m^v}T\chi}\le\exp(-m^s),
\]
because $\norm{\chi}\le1$.  Splitting at the previous scale gives
\begin{align*}
 \norm{T\chi}
 &\le \norm{\one_{H_A\le m^u}\chi}
       +\norm{T}\,\norm{\one_{H_A>m^u}\chi}\\
 &\le 1+\exp\bigl(m^k-m^{u-\alpha}\bigr)
 \le2.
\end{align*}
This is the first use of a global polynomial norm; it is paid for by
$k<u-\alpha$.

Proposition~\ref{prop:propagation}, applied with exponents $k<K$ to
the regular reference trajectory from $\chi$, gives an operator $E$
such that
\[
 \norm{VT\chi-EV\chi}\le\exp(-m^k),
 \qquad \deg E\le m^K,
 \qquad \norm{E}\le\exp(m^K).
\]
The endpoint $V\chi=P_BV\psi$ has exactly zero $H_B$ count.
By Lemma~\ref{lem:degree-shift},
$\one_{H_B>m^K}EV\chi=0$.  Therefore the filtered vector has the
effective-range bound
\[
 \norm{\one_{H_B>m^K}VT\chi}\le\exp(-m^k).
\]

\emph{Approximating the final projection on the controlled side.}
Use Lemma~\ref{lem:zero-filter} with precision exponent $q$, range
exponent $K$, and degree exponent $r$.  These satisfy $0<q<K<1$
and $r>(K+q)/2$.  Write $D=D(H_B)$ for the resulting self-adjoint
operator.  It has degree at most $m^r$, norm at most $\exp(m^r)$,
and error at most $\exp(-m^q)$ from $P_B$ on every integer count
at most $m^K$.  The preceding bound and $\norm{VT\chi}\le2$ imply
\begin{align*}
 \norm{(P_B-D)VT\chi}
 &\le 2\exp(-m^q)
       +\bigl(1+\exp(m^r)\bigr)\exp(-m^k)\\
 &=O\bigl(\exp(-m^q)\bigr).
\end{align*}
Here $q<r<k$ absorbs the global norm of $D$ on the remaining tail.
Only integer approximation is used, as required for the spectrum of
$H_B$.

Let $a_m=\ip{\chi}{T\chi}\ge0$.  Since $T$ is self-adjoint and
$V\chi=P_BV\psi$,
\[
 a_m=\ip{VT\chi}{P_BV\psi}.
\]
Moving the self-adjoint difference $P_B-D$ to the first argument
gives
\begin{align*}
 \abs{a_m-\ip{VT\chi}{DV\psi}}
 &=\abs{\ip{(P_B-D)VT\chi}{V\psi}}\\
 &\le\norm{(P_B-D)VT\chi}\,\norm{V\psi}
 =O\bigl(\exp(-m^q)\bigr).
\end{align*}
Thus the projection is approximated on $VT\chi$, the vector for which
an effective count range has just been proved.

\emph{Reverse propagation on the trajectory of $\psi$.}
The reversed trajectory starting at $V\psi$ is regular for
$V^\dagger$.  Applying Proposition~\ref{prop:propagation} to $D$
with exponents $r<r'$ gives an operator $G$ satisfying
\[
 \norm{V^\dagger DV\psi-G\psi}\le\exp(-m^r),
 \qquad \deg G\le m^{r'},
 \qquad \norm{G}\le\exp(m^{r'}).
\]
This application uses the unfiltered reference vector $V\psi$.
In the scalar product its error is multiplied by the vector norm
$\norm{T\chi}\le2$, so
\[
 \abs{a_m-\ip{T\chi}{G\psi}}
 \le O\bigl(\exp(-m^q)\bigr)+2\exp(-m^r)
 =O\bigl(\exp(-m^q)\bigr).
\]
The global operator norm of $T$ introduces no factor here.

\emph{Splitting the scalar product by the initial count.}
Put $Q=\one_{H_A\le m^v}$.  Its low-count contribution satisfies
\[
 \abs{\ip{QT\chi}{QG\psi}}
 \le\exp(-m^s)\,\exp(m^{r'}).
\]
For the high-count contribution, first observe that
$m^{v_0}+m^{r'}<m^v$ eventually.  The degree-shift property implies
\[
 (I-Q)G\one_{H_A\le m^{v_0}}=0.
\]
Using \eqref{eq:input-tail} with threshold exponent $v_0>b$
and precision exponent $\theta<v_0$ gives
\begin{align*}
 \norm{(I-Q)G\psi}
 &\le\norm{G}\,\norm{\one_{H_A>m^{v_0}}\psi}\\
 &\le\exp\bigl(m^{r'}-m^\theta\bigr).
\end{align*}
Orthogonality of $Q$ and $I-Q$ now yields
\begin{align*}
 \abs{\ip{T\chi}{G\psi}}
 &\le\exp\bigl(m^{r'}-m^s\bigr)
      +2\exp\bigl(m^{r'}-m^\theta\bigr)\\
 &=O\bigl(\exp(-m^q)\bigr),
\end{align*}
since $s>r'>q$ and $\theta>r'>q$.  These two strict gaps pay for
the two uses of the global norm of $G$.  In particular there is a
constant $C_0$, independent of $m$, with
\[
 0\le a_m\le C_0\exp(-m^q)
\]
eventually.  No commutation between $G$ and the count projections has
been assumed.

\emph{Positivity and descent.}
The tail filter is at least $1/4$ on $[m^h,m^u]$ and nonnegative
on the entire spectrum.  Consequently
\[
 \norm{\one_{m^h\le H_A\le m^u}\chi}^2
 \le4a_m\le4C_0\exp(-m^q).
\]
Combining this band with the previous upper tail gives
\[
 \norm{\one_{H_A\ge m^h}\chi}^2
 \le4C_0\exp(-m^q)
       +\exp\bigl(-2m^{u-\alpha}\bigr).
\]
Set $a=h-\beta>0$.  The exponent budget gives $q>a$, and
$u-\alpha-a=(u-h)+(\beta-\alpha)>0$.  Multiplying the last
right-hand side by $\exp(2m^a)$ therefore makes both terms tend to
zero.  Taking the square root proves
\[
 \norm{\one_{H_A\ge m^h}\chi}\le\exp(-m^{h-\beta})
\]
eventually, closing the descending step.  The strict power gap $q>a$
absorbs the square root and all fixed constants.

There are only $N$ descending steps.  Each uses
Proposition~\ref{prop:propagation} over the same bounded number of
layers with fixed positive exponent gaps $K-k$ and $r'-r$.
These gaps can be subdivided into that fixed number of positive
increments as in its proof.  Different descending steps retain only
the improved tail assertion; their approximating operators are not
composed.  Thus every choice and every number of steps is independent
of $m$, and the finitely many eventual estimates hold simultaneously
after increasing the threshold for $m$.

At the final scale,
\[
 h_N-\alpha_N=h_*-2\delta
 =\frac{h_*+p_*}{2}>p_*.
\]
The final induction bound implies the desired estimate at $h_*,p_*$.
Since these targets were arbitrary, the lemma follows.
\end{proof}
